\documentclass[dvipdfmx,a4j,twoside,openright]{jsbook}

\usepackage{newtxtext}
\usepackage{amsmath,amssymb}
\usepackage{mathtools}
%\usepackage{bm}

\usepackage[dvipdfmx]{graphicx}
%\usepackage{xcolor}
%\usepackage{tikz}

\usepackage{booktabs}
\usepackage{subcaption}

\usepackage{cite}
%\usepackage{pdfpages}

%\allowdisplaybreaks[1]
\numberwithin{equation}{chapter}

%\renewcommand{\bibname}{参考文献}

\newcommand{\Slash}[1]{{\ooalign{\hfil/\hfil\crcr$#1$}}}
%\usepackage{slashed}

%%
\DeclareMathOperator{\tr}{tr}
\DeclareMathOperator{\Tr}{Tr}
\DeclareMathOperator{\re}{Re}
\DeclareMathOperator{\im}{Im}
\DeclareMathOperator{\sign}{sign}

\usepackage[dvipdfm,margin=30truemm]{geometry}
\setlength{\textwidth}{150truemm}
\setlength{\fullwidth}{\textwidth}


%--- Don't add any more settings  ---%
%---  under this line in preamble ---%

% bookmarks of PDF with hyperref
\usepackage[dvipdfmx,%
bookmarks=true,bookmarksnumbered=true,%
pdftitle={},%
pdfsubject={Master thesis},%
pdfauthor={}%,
pdfkeywords={}%
]{hyperref}
\usepackage{pxjahyper}

%%%%%%%%%%%%%%%%%%%%%%%%%%%%%%%%%%%%%%%%%%%%%%%%%%%%%%%%%%%%%%%%
% TITLE
\title{何か面白い研究}
\author{九州大学~理学府~物理学専攻\\
粒子宇宙論講座~素粒子理論研究室\\
\\
Hogehoge~Foo\\%% author
\\\\
指導教員\hspace{3mm}
鈴木~博}
%\date{\today}
%\date{20xx年2月xx日}%回覧用
\date{20xx年2月xx日}%製本用
%%%%%%%%%%%%%%%%%%%%%%%%%%%%%%%%%%%%%%%%%%%%%%%%%%%%%%%%%%%%%%%%

\begin{document}
\maketitle
\cleardoublepage
\pagenumbering{roman}
%%%
\begin{abstract}
 \begin{center}
  \textgt{要旨}
 \end{center}
 
 要旨（丁寧に!!!）

「靴を脱げ。汝の立てるは聖なる地である。」 （\textgt{旧約聖書 出エジプト記}より）
\end{abstract}
%%%
\tableofcontents
\newpage
\pagenumbering{arabic}

%%%%%%%%%%%%%%%%%%%%%%%%%%%%%%%%%%%%%%%%%%%%%%%%%%%%%%%%%%%%%%%%
%本文
\chapter{序論}
とにかくintro, conclusionを丁寧に書こう!

``It's so fine and yet so terrible to stand in front of a blank canvas.'' (P. C\'ezanne)

\chapter{結論}
``No, no, you're not thinking; you're just being logical.'' (N. Bohr)
%%%%%%%%%%%%%%%%%%%%%%%%%%%%%%%%%%%%%%%%%%%%%%%%%%%%%%%%%%%%%%%%


%%%%%%%%%%%%%%%%%%%%%%%%%%%%%%%%%%%%%%%%%%%%%%%%%%%%%%%%%%%%%%%%
\chapter*{謝辞}
\addcontentsline{toc}{chapter}{謝辞}
感謝しなさい。

「この素晴らしい世界に祝福を!」（\textgt{この素晴らしい世界に祝福を!}より）
%%%%%%%%%%%%%%%%%%%%%%%%%%%%%%%%%%%%%%%%%%%%%%%%%%%%%%%%%%%%%%%%


%%%%%%%%%%%%%%%%%%%%%%%%%%%%%%%%%%%%%%%%%%%%%%%%%%%%%%%%%%%%%%%%
\appendix
\chapter{補足}
必要なら付録

``All work and no play makes Jack a dull boy.'' （\textit{Shining}より）
%%%%%%%%%%%%%%%%%%%%%%%%%%%%%%%%%%%%%%%%%%%%%%%%%%%%%%%%%%%%%%%%


%%%%%%%%%%%%%%%%%%%%%%%%%%%%%%%%%%%%%%%%%%%%%%%%%%%%%%%%%%%%%%%%
%"適度な"文献
\begin{thebibliography}{00}
%\cite{Ishikawa:2019tnw}
\bibitem{Ishikawa:2019tnw}
  K.~Ishikawa, O.~Morikawa, A.~Nakayama, K.~Shibata, H.~Suzuki and H.~Takaura,
  %``Infrared renormalon in the supersymmetric $\mathbb{C}P^{N-1}$ model on $\mathbb{R}\times S^1$,''
  arXiv:1908.00373 [hep-th].
  %%CITATION = ARXIV:1908.00373;%%
\end{thebibliography}

%BibTeX使うなら（ref.bibに文献記入）
%\bibliographystyle{utphys}
%\bibliography{ref}
%%%%%%%%%%%%%%%%%%%%%%%%%%%%%%%%%%%%%%%%%%%%%%%%%%%%%%%%%%%%%%%%

\end{document}
